\properlane{household-work}{The work the master leaves in the house}

\subsection{What has the owner put into his servants' hands?}

A householder about to go away must decide what can safely be left until he gets back. In Jesus' story he is confident enough in the servants to entrust them with authority. Yet he issues a command to keep watch. What does he expect to find them doing on his return? A porter motionless beside a closed door might be watching as required, or he might have fallen asleep. The other servants might be looking down the road with such eagerness for their master's return that they have left their own jobs unattended. Jesus describes a house in which the owner expects work to go on while he is away. His hearers therefore have two things to consider, the man they wish to see again and the use they make of what he left in their hands.

Bede asks what kind of journey Christ made when he left his Church. Christ rose and ascended bodily to the Father. The flesh proper to life here on earth has been taken into heaven. Yet the risen Christ has also promised to remain with his people. A person setting off on an ordinary journey loses the power to help a servant at the far end of the house. What Christ does for his servants in the course of his absence exceeds anything we might infer simply from the distance the traveler has gone. According to Bede, he gives his faithful the grace of the Holy Spirit. That is how he gives them the ability to set about the good works at their disposal. Christ can be absent in the body which ascended and present as the one who enables the work to go forward.\footnote{Bede, \textit{In Marci evangelium expositio}, on Mark 13:33--37, Latin text, lines 403--405, especially the comments on the master's journey and his gifts at line 404.}

The solemnity of the command to the porter makes sense if the people inside the house stand to be harmed by what comes through the door. Bede first thinks of the responsibility of pastors. They owe attention to people committed to their protection. Neglect at the door can endanger someone besides the man responsible for it. But Jesus ends by extending to \emph{all} what he has just said to his disciples. That final word prevents Bede from leaving the entire obligation with the successors of the apostles. There is an entrance for which each person must answer, the entrance into the heart through which an evil suggestion can gain a hearing. Where one servant has charge of a community and another has no such office, neither can hire someone to watch over his own consent to evil. Each has some power which the other cannot exercise for him.

For Bede, acting upon one's belief belongs to staying awake. A servant might believe the owner will return and still leave the door unguarded night after night. His belief about the future may be correct. Everything depends, for the people who live behind that door tonight, on whether he will keep his post. Belief without the exercise of the charge leaves them as exposed as if he had denied that there was any owner. There is also a place for acts nobody else can observe. A pastor might publicly perform his duties while secretly turning another person's confidence to his own advantage. For that danger he stands in need of the same custody of the heart as any other Christian. No external office makes the final charge superfluous.

\subsection{How can a church lack no gift and still be in danger?}

The Corinthians offer an example of people abundantly supplied for the work of a house. Paul speaks of them being enriched in utterance and knowledge. He has grounds for his thanksgiving. Yet later he has to address them as people little advanced in spiritual life (1 Cor 3:1). Chrysostom puts the question bluntly. How could the apostle, who knew the state they were in well enough to rebuke them, tell these same people they lacked nothing? Telling a divided house how splendidly it is doing can only deepen its self-deception. At the same time Paul has seen things in their midst for which he owes thanks to God.

Chrysostom first distinguishes a congregation's beginnings from its present life. People who once received their gifts eagerly may subsequently have grown slack. Praising what they had received brings back the memory of those first days precisely when remembering might make them ashamed of how they now behave. There is also a difference between the members who deserve a rebuke and those who go on serving. In the later parts of the letter one person still has a psalm, another something to teach, another a gift of interpretation. Paul can address the same church both ways because it contains people for whom his assessments have different force. Chrysostom also allows for ordinary collective speech. Paul can say of the whole church what holds for the greater part. Why open so gently in any of these cases? A weaker hearer might take offense at the first reproach and never hear the correction. What sounds like praise of the Corinthians gives the credit for their forgiveness and their new standing before God to his grace. Paul reminds people who must face their faults how generously God has already treated them.\footnote{John Chrysostom, \textit{Homilies on First Corinthians}, Homily 2, §§4--5, on 1 Cor 1:7 and his comparisons with 3:1 and 14:26--29.}

Thomas Aquinas asks who possesses the gifts for which Paul gives thanks. In verse 5 Paul speaks to the Corinthians as people enriched in speech and knowledge. Thomas cites Augustine on the person who loves the unity of the Church. Envy makes another Christian's ability a cause for resentment. Charity allows a member who lacks that ability to welcome it as a good that already belongs to him through his neighbor. In verse 7 Thomas explains how a church with many members can lack no gift. God has distributed the gifts among different people. A learner listening to a gifted speaker illustrates the distinction. He may learn something he could never have explained on his own. But even before the lesson, he can rejoice that someone in the church they both love has the ability to teach. His share in that good begins with charity and requires no talent for speaking on his part.\footnote{Thomas Aquinas, \textit{Super I ad Corinthios}, chapter 1, lecture 1, §86234, on 1 Cor 1:5--7, citing Augustine on John at v.~5.}

Paul ends his thanksgiving by naming the fellowship to which the Corinthians belong. A few verses later one man claims Paul and another Apollos. There is a cost to treating a preacher who was sent to the whole church as the possession of a party. Those who claim him narrow the circle of people to whom his gift is available. Those outside the party find themselves tempted to refuse something from which they too could benefit. Paul's famous question about who was crucified for them tests these loyalties at their root. Can the teacher claimed by one group be the source of the life given to the others? The congregation already heard in the greeting that grace and peace come from the Father and the Lord Jesus Christ. The good thing a teacher can do for them depends on a relation already established by another.

When Paul speaks of the day on which Christ will be revealed, Chrysostom takes him to mean that the one who will come is present even though unseen. Being revealed involves becoming visible. Then Paul promises that the Corinthians will be strengthened to the end. A promised strengthening tells them something about their present need, as a hand offered to a person on uncertain ground tells him something about his footing. They require help even after an abundance of gifts has been given them. It is possible to learn that from a promise without losing courage. God's fidelity allows them to take the danger seriously enough to submit to correction. Judged on their immediate prospects of repairing their disputes unaided, they would have good reason to despair.\footnote{Chrysostom, \textit{Homilies on First Corinthians}, Homily 2, §§6--8.}

\subsection{Which servants have time to do this work?}

If Jesus spoke the closing word of the Gospel to everyone, what of the people who heard it centuries ago and have already died? In his letter to Hesychius, Augustine mentions Mark's list of watches and then brings the question to bear on the people he himself knows. The first hearers have gone. There will be Christians after the people to whom Augustine writes. He refuses to assign the command only to those who happen to be alive at the end of the world. When people die, the manner in which they lived and died will matter when they stand before the judge. There is therefore something at stake for any generation that hears the command today, even if it will have no occasion to see the Lord's public coming in its earthly lifetime. A person does not have the interval between his own death and the general judgment in which to do the work neglected during life. Augustine preserves both a real end to each person's time for acting and the judgment still awaited by the whole Church.\footnote{Augustine, \textit{Letter} 199 to Hesychius, I.2--3, especially his citation of Mark in I.3.}

Thus Bede and Augustine can attend to Jesus' word \emph{all} along different lines. Bede is concerned with someone who may think the porter or the pastor has taken charge of everything. Augustine is concerned with someone who might think that another age, nearer the end of time, will have a better reason to take the warning seriously. Neither excuse is available to a Corinthian who has the chance to lay aside his partisanship today. Nor is it available to a Christian who has just prayed the Collect. The Christian has asked for the will to go to meet Christ with deeds of justice. For someone called to a particular piece of work, justice may begin with making it possible for one of the other servants to do his.

In Isaiah the congregation wishes that God might encounter it doing right, then confesses that nobody troubled to take hold of him. There has already been an occasion to exercise the fidelity the people now desire. In the psalm, similarly, the vine has an established place in the earth. Its shoots and branches depend upon the hand that planted it. The people of Israel ask God to tend their own vine, whose history they trace back to Egypt. Augustine and Bellarmine both turn to Paul's olive tree in Romans 11 to explain how the Gentiles came to share in that history. Augustine dwells on the request that God perfect what his right hand planted. The vine is Abraham's seed, to whom God promised that all nations would be blessed. The wild olive branch depends on the root into which it has been grafted. Augustine identifies Christ as the foundation on which that original planting reaches its fulfillment. Bellarmine asks who actually lived through the vineyard's reform. The apostles were Israelites, as were the thousands converted through Peter's preaching. When Gentiles believed, they too were grafted onto Israel's tree. Only then does Bellarmine speak of a better Church in the old vineyard's place. His emphasis falls on the charge given to the Son of Man. Christ is the chief husbandman under whose care the vineyard can yield the fruit for which God planted it.\footnote{Augustine, \textit{Exposition on Psalm} 79, §9, citing Gen 22:18 and 1 Cor 3:11. Robert Bellarmine, \textit{Commentary on the Book of Psalms}, Psalm 79, explanations 14--15.}

\subsection{What the household receives from its Lord}

At the Alleluia the people's petition for salvation can be heard as a request for the Savior in person. Bellarmine takes it this way as he explains the psalm. At Communion he attends to a second verse in that psalm which speaks of fruit. Much of a farmer's work consists in preparing the soil, sowing seed, cultivating what grows and keeping animals out until harvest. Bellarmine pictures a different kind of growth. A flower of the field comes up without human sowing, under the sun and rain. He understands the land as the Virgin and God's gift from above as the work of the Holy Spirit. The fruit is Christ born of Mary. At Mass those who approach the altar have already brought the products of work actually done in fields and vineyards. What Bellarmine says of the fruit allows us to ask who first made that human offering possible. The one given there is the Son, whose coming in the flesh no cultivation by the servants could have produced.\footnote{Bellarmine, \textit{Commentary on the Book of Psalms}, Psalm 84, explanations 7 and 12. Cyril of Jerusalem, \textit{Catechetical Lecture} XV.1--3, distinguishes the humility of the first coming from the majesty of the second.}

The U.S. commentary on the prayer over the gifts also mentions what people produce to feed the hungry. No one servant can provide for the household simply by being willing to work. Someone who cultivates the field, someone who makes bread and someone in a position to give it away can each do something for the others. After Communion the people ask about the proper use of passing things. Bread does perish, and so does an opportunity to feed somebody with it. The act has a lasting bearing on the relation of the servants to their master and to each other. This is one reason to ask for instruction in the use of particular goods as they come into one's hands. The Advent blessing, when pronounced, commends precisely such a present exercise of charity to those awaiting the return. The hope expressed at the Entrance has companions already in the house. The person who lifts his soul to God can look to the people who have waited before him and count on the fidelity by which God sustained them.

\begin{fourSenses}
\item[Literal.] The master distributes work and charges the porter to watch. Jesus addresses the final warning to all. Paul can thank God for the endowments of a church whose members still need to learn to serve each other with them.
\item[Allegorical.] In Bede's account the master is Christ bodily ascended to the Father. The Spirit he gives makes the work of his Church possible during his absence. The humbled Son whose first coming the Church recalls is the judge she expects to see.
\item[Moral.] A servant is answerable for the duties of his office and for what he admits into his heart. Even a servant who lacks another's ability can put aside envy and rejoice in the good it brings to the church they both love. Through charity he already shares in his neighbor's gift.
\item[Anagogical.] Every generation of Christians will answer to the returning Lord. A servant who dies before the final day still awaits judgment for a life whose time for service has run its course.
\end{fourSenses}
